% ---------------------------------------------------------------
\section{Results and Analysis}
\label{sec:results}

% generated by scripts/build_tables.mjs from results/phase3/v2_final/impact_summary.csv, results/phase3/v2_final/harm_counts.csv, results/phase3/v2_final/input_detection_summary.csv, results/phase3/v2_final/impact_per_seed.csv, results/phase3/v2_final/paired_estimator_bootstrap.csv
\begin{table*}[t]
\centering
\caption{Controlled synthetic evaluation, logistic monitor, 30 seeds $\times$ 24 batches per condition (all batches). Detect.: seeds whose onset (batches 9--12) was matched by an input alarm, defined only for input-changing conditions. Alarm frac.: fraction of post-onset batches with an input alarm (all 24 batches for C0). MAE: mean absolute error of the batch-accuracy estimate, with $\pm$ the standard deviation over seeds for the calibrated estimator; lowest per row in bold. $^{\dagger}$/$^{\ddagger}$: the comparator's MAE is significantly higher/lower than the calibrated estimator's (paired seed-bootstrap 95\% interval of the difference excludes 0). Cov.: empirical coverage of the nominal 95\% interval ($\pm$ SD). Harm: harmful batches of 720. TP/FP: concern flags ($\kappa_t\ge0.95$) on harmful / non-harmful batches.}
\label{tab:core}
\setlength{\tabcolsep}{3.5pt}
\begin{tabular}{llcccccccrc}
\toprule
 & & \multicolumn{2}{c}{Input channel} & \multicolumn{4}{c}{Accuracy-estimate MAE} & & \multicolumn{2}{c}{Harm concern} \\
\cmidrule(lr){3-4}\cmidrule(lr){5-8}\cmidrule(lr){10-11}
Cond. & Change & Detect. & Alarm frac. & Calibrated AC & Raw AC & ATC-style & Constant & Cov. & Harm & TP/FP \\
\midrule
C0 & none (stationary) & -- & 0.05 & 0.017\,{\scriptsize$\pm$0.004} & \textbf{0.016} & 0.023$^{\dagger}$ & 0.019 & 0.92\,{\scriptsize$\pm$0.07} & 11 & 0/0 \\
C1 & mean, irrelevant $x_7$ & 30/30 & 1.00 & 0.017\,{\scriptsize$\pm$0.004} & \textbf{0.016} & 0.023$^{\dagger}$ & 0.019 & 0.92\,{\scriptsize$\pm$0.08} & 11 & 0/0 \\
C2 & gradual mean, $x_0$ & 30/30 & 0.97 & 0.016\,{\scriptsize$\pm$0.004} & \textbf{0.015}$^{\ddagger}$ & 0.039$^{\dagger}$ & 0.063$^{\dagger}$ & 0.90\,{\scriptsize$\pm$0.08} & 6 & 0/0 \\
C3 & $P(Y|X)$ reversal & -- & 0.05 & \textbf{0.303}\,{\scriptsize$\pm$0.007} & 0.303 & 0.304 & 0.303 & 0.31\,{\scriptsize$\pm$0.04} & 485 & 0/0 \\
C4 & gradual reversal & -- & 0.05 & \textbf{0.266}\,{\scriptsize$\pm$0.007} & 0.267 & 0.267 & 0.266 & 0.31\,{\scriptsize$\pm$0.04} & 473 & 0/0 \\
C5 & recurrent reversal & -- & 0.05 & 0.160\,{\scriptsize$\pm$0.005} & \textbf{0.159} & 0.164$^{\dagger}$ & 0.161 & 0.61\,{\scriptsize$\pm$0.05} & 247 & 0/0 \\
C6 & corr$(x_0,x_1)$ .3$\to$.9 & 8/30 & 0.13 & 0.019\,{\scriptsize$\pm$0.004} & \textbf{0.017}$^{\ddagger}$ & 0.074$^{\dagger}$ & 0.057$^{\dagger}$ & 0.90\,{\scriptsize$\pm$0.09} & 414 & 160/9 \\
\bottomrule
\end{tabular}
\end{table*}


\subsection{Input Evidence Separates from Harm}
\label{sec:res-core}
Table~\ref{tab:core} reports the core outcome for the logistic monitor, and
Fig.~\ref{fig:heatmap} condenses five indicators for all twelve synthetic conditions
into one evidence map. Marginal changes are found reliably. The irrelevant shift C1
and the gradual relevant shift C2 were both matched in 30/30 seeds; C2's mean
first-alarm delay was 0.43 batches (95\% interval 0.27--0.63). On the stationary stream
the Holm-corrected KS family raised 48.6 alarms per 1,000 batches [29.2, 73.6],
consistent with its familywise level of 0.05. Detection, however, does not imply harm.
C1 alters only an unused feature: it raised alarms in every post-onset batch, yet only
11 of its 720 batches met the harm rule~\eqref{eq:harm}, as many as the stationary C0.

\begin{figure}[t]
\centering
% Evidence map. Data: data/heatmap.dat <- results/phase3/v2_final/{input_detection_summary,harm_counts,impact_summary}.csv
% (logistic monitor, regime all, 30 seeds x 24 batches per row). Generated by scripts/build_tables.mjs.
\begin{tikzpicture}
\begin{axis}[
  idm, axis lines=box, axis line style={draw=none}, ymajorgrids=false, clip=false,
  width=0.80\columnwidth, height=8.2cm,
  enlargelimits=false, xmin=-0.5, xmax=4.5, ymin=-0.5, ymax=11.5,
  xtick={0,...,4}, xticklabels={Input\\alarm, Harm\\rate, Concern\\recall, Interval\\coverage, Estimate\\MAE},
  x tick label style={align=center, font=\scriptsize\bfseries, text=cink}, xticklabel pos=upper,
  ytick={0,...,11}, yticklabels={X4,X3,X2,X1,X0,C6,C5,C4,C3,C2,C1,C0},
  y tick label style={font=\scriptsize\bfseries, text=cink},
  tick style={draw=none},
  colormap={idmheat}{color=(white) color=(c6!45) color=(c1!80)},
  point meta min=0, point meta max=1,
  colorbar horizontal,
  colorbar style={height=2.2mm, at={(0.5,-0.03)}, anchor=north, xtick={0,0.5,1},
    every tick label/.append style={font=\tiny}, xlabel={cell shade (MAE column scaled by 1/0.35)}, xlabel style={font=\tiny, yshift=4pt}},
]
\addplot[matrix plot*, mesh/cols=5, point meta=\thisrow{m}, draw=white, line width=1.2pt,
  visualization depends on={\thisrow{v} \as \rawv},
  nodes near coords={\pgfmathprintnumber[fixed, fixed zerofill, precision=2]{\rawv}},
  every node near coord/.append style={font=\scriptsize, text=cink, anchor=center},
] table[x=x, y=y] {data/heatmap.dat};
\draw[cink, line width=0.9pt] (axis cs:-0.5,4.5) -- (axis cs:4.5,4.5);
\draw[c2, line width=1.2pt, rounded corners=1.5pt] (axis cs:-0.47,5.53) rectangle (axis cs:4.47,8.47);
\node[font=\scriptsize\itshape, text=c2, anchor=west, align=left] at (axis cs:4.55,7) {silent\\failure};
\node[font=\scriptsize\itshape, text=cgray, anchor=west, align=left] at (axis cs:4.55,8.6) {linear\\task};
\node[font=\scriptsize\itshape, text=cgray, anchor=west, align=left] at (axis cs:4.55,2) {XOR\\task};
\end{axis}
\end{tikzpicture}

\caption{Evidence map of the controlled benchmark (logistic monitor; 30 seeds $\times$ 24 batches per row). Input alarm: fraction of post-onset batches with an alarm. Harm rate: harmful batches / 720. Concern recall: harmful batches flagged by $\kappa_t\ge0.95$. Coverage: empirical coverage of the nominal 95\% interval. Estimate MAE: mean absolute error of $\hat A_t$ (its shading is scaled by $1/0.35$). The outlined rows (C3--C5) show high harm with input alarms at the false-alarm rate and zero recall.}
\label{fig:heatmap}
\end{figure}

The opposite failure is larger, and it stands out in the outlined rows of
Fig.~\ref{fig:heatmap}. Under C3--C5 the input alarm fired in 5\% of post-onset
batches, the stationary rate, because the inputs do not change. Accuracy nevertheless
collapsed: 485, 473 and 247 batches were harmful, and the calibrated concern rule
flagged none of them, for either learner. The dependence-only change C6 lies between
the two cases. It was matched at onset in only 8/30 seeds, and the concern flag
recovered 160 of its 414 harmful batches at the cost of 9 false flags. Across all
conditions the logistic monitor left on average 92.8\% of each seed's harmful batches
unflagged (recall 0.072 [0.043, 0.103]; 90.3\% pooled over batches), with a false-flag
rate of 0.004. The monitor is therefore conservative, but it misses most harm.

Estimation accuracy follows the same split (Table~\ref{tab:core}). When only $P(X)$
changes (C0--C2 and C6), the calibrated AC tracks batch accuracy to an MAE of
0.016--0.019 (seed SD about 0.004), and its 95\% interval covers 90--92\% of batches.
It is significantly more accurate than the ATC-style estimator in every one of these
conditions, and more accurate than the constant reference wherever accuracy moves (C2,
C6). Raw confidence is slightly but significantly better on C2 and C6, because the
logistic model is already close to calibrated. Under C3 the MAE rises to 0.303 and
coverage falls to 0.31. No unlabelled estimator escapes this failure: the
differences between estimators on C3 and C4 are not significant.

\subsection{Conditional-Only Change: a Silent Failure}
\label{sec:res-conditional}
Fig.~\ref{fig:fail} averages 30 seeds per batch. Before onset, estimate and realised
accuracy agree. From batch~9 the estimate stays at its pre-onset level, while realised
accuracy falls to between 0.27 and 0.28 once a reversal is complete. C4 reaches this
level at batch 13, and accuracy recovers only when C5 restores the law at batch 17.
Over batches 9--24 of C3--C5 the mean estimate is 0.723 and the mean realised accuracy
is 0.370. The classifier receives familiar inputs and so reports familiar confidence.
A calibrator fitted on the source law cannot tell that the law has reversed.

\begin{figure}[t]
\centering
\resizebox{\columnwidth}{!}{% Data: data/conditional_failure.dat <- reports/phase3/v2_final/v1_conditional_failure_diagnostics.csv
% (mean over 30 seeds per batch, logistic monitor). Built by scripts/build_tables.mjs.
\begin{tikzpicture}
\begin{axis}[
  idm, width=\columnwidth, height=5.0cm,
  xmin=0.5, xmax=24.5, ymin=0.2, ymax=0.97,
  xlabel={batch $t$ (500 rows each)}, ylabel={batch accuracy},
  xtick={1,5,9,13,17,21,24},
  legend style={at={(0.5,-0.32)}, anchor=north, legend columns=2},
]
\fill[c2!7] (axis cs:8.5,0.2) rectangle (axis cs:24.5,0.97);
\draw[c2!60, dashed, thick] (axis cs:8.5,0.2) -- (axis cs:8.5,0.97);
\node[font=\scriptsize\bfseries, text=c2, anchor=north west] at (axis cs:8.7,0.96) {onset: $P(Y\mid X)$ changes, $P(X)$ does not};
\addplot[c1, line width=1.6pt] table[x=batch, y=est] {data/conditional_failure.dat};
\addlegendentry{estimate $\hat A_t$ (same for C3--C5)}
\addplot[c2, thick, mark=*, mark size=1.1pt] table[x=batch, y=c3] {data/conditional_failure.dat};
\addlegendentry{realised $A_t$: C3 abrupt}
\addplot[c3, thick, mark=square*, mark size=1.0pt] table[x=batch, y=c4] {data/conditional_failure.dat};
\addlegendentry{realised $A_t$: C4 gradual}
\addplot[c4, thick, mark=triangle*, mark size=1.4pt] table[x=batch, y=c5] {data/conditional_failure.dat};
\addlegendentry{realised $A_t$: C5 recurrent}
\draw[{Stealth[length=1.8mm]}-{Stealth[length=1.8mm]}, cink!70, thick] (axis cs:15,0.2813) -- (axis cs:15,0.7189) node[midway, right, font=\scriptsize, text=cink, fill=white, inner sep=1pt, align=left] {silent gap 0.44\\(batch 15)};
\end{axis}
\end{tikzpicture}
}
\caption{Calibrated accuracy estimate versus realised accuracy under conditional-only change (logistic monitor, mean over 30 seeds per batch). The estimate is identical across C3--C5 because their inputs are identical. The shaded region is post-onset.}
\label{fig:fail}
\end{figure}

Fig.~\ref{fig:case} shows the same failure as an analyst would see it in the aligned
view, for one stream (seed 100, C3, delay 5) at clock 16. For batch 10 the estimate is
0.731 [0.694, 0.768], the concern score is 0.0006, and the smallest Holm-adjusted
$p$-value is 0.174, so no alarm fires. The released accuracy of batch 10, available
from clock 15, is 0.254. Batches 12--16 are still unresolved and are shown as such.
The only input alarm in this stream occurred at batch 6, before onset, and was false.
The display exposes the disagreement once labels are released; before then nothing in
the unlabelled evidence signals concern.

\begin{figure}[t]
\centering
\resizebox{\columnwidth}{!}{% Qualitative example (vector redraw of the evidence-replay view state: seed 100, C3, logistic, delay 5, as-of clock 16).
% Data: data/case_c3.dat <- results/phase3/v2_final/predictions/v1_synthetic_100_C3_logistic_online.csv
%       and results/phase3/v2_final/released_verification.csv. Values cross-checked against the HTML replay view.
\begin{tikzpicture}
\begin{axis}[
  idm, width=\columnwidth, height=5.0cm,
  xmin=0.5, xmax=16.5, ymin=0.15, ymax=0.9,
  xlabel={batch $b$ (view as of clock $c=16$, delay $d=5$)}, ylabel={batch accuracy},
  xtick={1,...,16}, unbounded coords=jump,
  legend style={at={(0.5,-0.34)}, anchor=north, legend columns=2},
]
\fill[cink!8] (axis cs:11.5,0.15) rectangle (axis cs:16.5,0.9);
\node[font=\scriptsize, align=center, text=cink!80] at (axis cs:14,0.45) {labels not yet\\released:\\\textbf{unresolved}};
\draw[c2!60, dashed, thick] (axis cs:8.5,0.15) -- (axis cs:8.5,0.9);
\node[font=\scriptsize\bfseries, text=c2, anchor=north west] at (axis cs:8.6,0.89) {onset};
\addplot[name path=lo, draw=none, forget plot] table[x=batch, y=lo] {data/case_c3.dat};
\addplot[name path=hi, draw=none, forget plot] table[x=batch, y=hi] {data/case_c3.dat};
\addplot[c1!18, forget plot] fill between[of=lo and hi];
\addplot[c1, line width=1.5pt, mark=*, mark size=1pt] table[x=batch, y=est] {data/case_c3.dat};
\addlegendentry{estimate + 95\% band ($\kappa_t<0.01$)}
\addplot[c2, thick, mark=square*, mark size=1.3pt] table[x=batch, y=acc] {data/case_c3.dat};
\addlegendentry{released accuracy ($b+d\le c$)}
\addplot[only marks, mark=triangle*, mark size=2.6pt, c3] coordinates {(6,0.81)};
\addlegendentry{input alarm (pre-onset, false)}
\node[draw=c2, fill=white, rounded corners=1pt, font=\scriptsize, align=left, inner sep=2pt, anchor=south west] at (axis cs:1.2,0.2) {batch 10: est.\ 0.731, released 0.254;\\no input alarm, no concern flag};
\end{axis}
\end{tikzpicture}
}
\caption{Qualitative example from the aligned evidence view: one replayed stream (seed 100, C3, logistic monitor, delay 5) as of clock 16. Accuracy appears only for batches whose labels are released ($b+5\le16$). Batches 12--16 are unresolved and are left blank rather than imputed.}
\label{fig:case}
\end{figure}

\subsection{Localization}
When the changed feature is known, the KS ranking localizes it. Over 30 seeds, P@1 and
R@3 were 1.0 and 1.0 for C1 and 0.988 and 0.994 for C2. On the 96-batch tasks the
correlation-change pair score alarmed in every active dependence window of the linear
and XOR tasks and named the programmed pair every time. On the five-region task its
alarm rate was 0.955 [0.937, 0.970] and its pair accuracy 0.991, with a stationary
alarm rate of 0.054--0.066 per batch. Frozen clusters did not recover the generating
regions: with five clusters the adjusted Rand index was 0.088 [0.080, 0.097], and
recall for the intervened region was 0.216, close to its prevalence of about 20\%.

\subsection{Natural Streams}
\label{sec:res-natural}
% generated by scripts/build_tables.mjs from results/phase3/uci224_delayed_replay_20261116_v2/evaluation/summary.csv, results/phase3/external_uci224_exact_interval_20261116_v1/summary.json, results/phase3/elec2_calibration_family_20261116_v1/summary.csv, results/phase3/elec2_temporal_20261116_v2/evaluation/summary.json, results/phase3/noaa_weather_1949_1999_v1/run_summary.json, results/phase3/noaa_weather_1949_1999_v1/evaluator/batch_metrics.csv
\begin{table*}[t]
\centering
\caption{Accuracy-estimation error on three natural streams (equal-batch mean absolute error, MAE). ELEC2 and NOAA use the isotonic correctness map; UCI224 uses a one-dimensional logistic correctness map. \emph{Cumulative} recalibrators refit on labels released strictly before each forecast. ECE: 10-bin expected calibration error. Cov.: batches whose 95\% interval contained the realised accuracy.}
\label{tab:external}
\setlength{\tabcolsep}{3.5pt}
\begin{tabular}{llccc}
\toprule
Stream & Estimator (delay) & MAE & ECE & Cov. \\
\midrule
\multirow{4}{*}{\shortstack[l]{UCI224\\9 cohorts}} & Frozen calibrated AC & 0.545 & 0.545 & 0/9 \\
 & Raw AC & 0.411 & 0.412 & -- \\
 & Historical constant & 0.512 & 0.512 & -- \\
 & Cumulative (0 / 1 / 2) & 0.169 / 0.227 / 0.283 & 0.188 / 0.245 / 0.290 & -- \\
\midrule
\multirow{3}{*}{\shortstack[l]{ELEC2\\55 batches}} & Frozen isotonic & 0.059 & 0.080 & 23/55 \\
 & Cum.\ isotonic (0/5/20) & 0.063 / 0.064 / 0.059 & 0.086 / 0.089 / 0.082 & 18/21/24 \\
 & Cum.\ Platt (0/5/20) & 0.065 / 0.067 / 0.062 & 0.102 / 0.105 / 0.103 & 21/21/24 \\
\midrule
\multirow{4}{*}{\shortstack[l]{NOAA\\21 batches}} & Frozen calibrated AC & 0.044 & -- & 9/21 \\
 & Raw AC & 0.088 & -- & -- \\
 & ATC-style & 0.067 & -- & -- \\
 & Historical constant & 0.044 & -- & -- \\
\bottomrule
\end{tabular}
\end{table*}


Table~\ref{tab:external} and Figs.~\ref{fig:uci}--\ref{fig:elec2} report the natural
streams. The UCI gas sensor stream is the hardest case. The frozen estimate stayed
between 0.975 and 0.977 for every cohort while realised accuracy ranged from 0.194 to
0.658; all nine cohorts raised input alarms, and the 95\% interval covered none. The
exact Poisson-binomial interval was slightly narrower than its normal approximation
(mean width 0.0234 against 0.0258), so approximation error does not explain the
failure: calibration did not transfer to the drifted sensors. Recalibrating the
correctness map on labels released before each forecast reduced MAE from 0.545 to
0.169, 0.227 and 0.283 at delays of 0, 1 and 2 cohorts. The benefit shrinks with delay
because each forecast is fitted on older cohorts, and the estimates remain biased
upward (Fig.~\ref{fig:uci}). Split and learner choices matter as much as delay: over
six stratified splits, the calibrated MAE on cohorts 3--10 was 0.553 (logistic) and
0.374 (tree) with cohort~1 as history, and 0.342 and 0.401 with cohorts 1--2.

\begin{figure}[t]
\centering
\resizebox{\columnwidth}{!}{% Data: data/uci_cohorts.dat <- results/phase3/external_uci224_day2_v2/cohort_metrics.csv (accuracy, estimated_accuracy)
% and results/phase3/uci224_delayed_replay_20261116_v2/evaluation/per_batch_metrics.csv (cumulative_released_logistic).
\begin{tikzpicture}
\begin{axis}[
  idm, width=\columnwidth, height=5.0cm,
  xmin=1.7, xmax=10.3, ymin=0.1, ymax=1.05,
  xlabel={UCI gas-sensor acquisition cohort}, ylabel={cohort accuracy},
  xtick={2,...,10},
  legend style={at={(0.5,-0.32)}, anchor=north, legend columns=3},
]
\addplot[name path=est, c1, line width=1.5pt, dashed] table[x=batch, y=frozen] {data/uci_cohorts.dat};
\addlegendentry{frozen estimate}
\addplot[name path=acc, cink, line width=1.6pt, mark=*, mark size=1.5pt] table[x=batch, y=acc] {data/uci_cohorts.dat};
\addlegendentry{realised accuracy}
\addplot[c2!12, forget plot] fill between[of=est and acc];
\addplot[c2, thick, mark=square*, mark size=1.1pt] table[x=batch, y=d0] {data/uci_cohorts.dat};
\addlegendentry{recalibrated, $d{=}0$}
\addplot[c3, thick, mark=triangle*, mark size=1.4pt] table[x=batch, y=d1] {data/uci_cohorts.dat};
\addlegendentry{recalibrated, $d{=}1$}
\addplot[c4, thick, mark=diamond*, mark size=1.4pt] table[x=batch, y=d2] {data/uci_cohorts.dat};
\addlegendentry{recalibrated, $d{=}2$}
\node[font=\scriptsize, text=c2, align=center] at (axis cs:7.2,0.83) {frozen-estimate error\\(MAE 0.545)};
\end{axis}
\end{tikzpicture}
}
\caption{UCI gas sensor stream: realised cohort accuracy, the frozen calibrated estimate (the shaded gap is its error), and estimates recalibrated on labels released at delays of 0, 1 and 2 cohorts.}
\label{fig:uci}
\end{figure}

\begin{figure}[t]
\centering
\resizebox{\columnwidth}{!}{% Data: data/elec2.dat <- results/phase3/elec2_temporal_20261116_v2/evaluation/evaluator_batches.csv (frozen monitor).
\begin{tikzpicture}
\begin{axis}[
  idm, width=\columnwidth, height=4.5cm,
  xmin=0.5, xmax=55.5, ymin=0.4, ymax=0.95,
  xlabel={ELEC2 monitoring batch (500 rows; batch 55 has 188)}, ylabel={batch accuracy},
  legend style={at={(0.99,0.98)}, anchor=north east, legend columns=2},
]
\addplot[name path=lo, draw=none, forget plot] table[x=batch, y=lo] {data/elec2.dat};
\addplot[name path=hi, draw=none, forget plot] table[x=batch, y=hi] {data/elec2.dat};
\addplot[c1!18, forget plot] fill between[of=lo and hi];
\addplot[c1, line width=1.5pt] table[x=batch, y=est] {data/elec2.dat};
\addlegendentry{estimate + 95\% band}
\addplot[c2, thick, mark=*, mark size=0.9pt] table[x=batch, y=acc] {data/elec2.dat};
\addlegendentry{realised (after release)}
\end{axis}
\end{tikzpicture}
}
\caption{ELEC2: frozen calibrated estimate with its 95\% plug-in band, and realised batch accuracy, available only after release. Input alarms fired in all 55 batches (not shown).}
\label{fig:elec2}
\end{figure}

ELEC2 shows the opposite pattern. Accuracy was 0.628 and the frozen estimate erred by
0.059 on average. Input alarms fired in all 55 batches, so an alarm carried no
information about which batches were harmful. Cumulative isotonic and Platt
recalibration removed most of the signed bias (from $-0.0106$ to between $-0.0005$ and
$-0.0023$) but \emph{increased} MAE at delays 0 and 5 (isotonic: 0.063 and 0.064) and
left it essentially unchanged at delay 20 (0.059); ECE did not improve at any delay.
On NOAA the calibrated estimate (MAE 0.044) matched the constant historical estimate
(0.044) and improved on raw confidence (0.088) and on the ATC-style estimate (0.067);
its interval covered only 9/21 batches, and the input alarm fired in every batch.
Selective-label acquisition on NOAA is analysed in Appendix~\ref{app:extra}.

\subsection{Prospective Harm Forecasting}
\label{sec:res-forecast}
A separate study tested whether harm over the next five batches can be predicted.
Forecasters were trained on development seeds and evaluated on held-out seeds, using
only labels released strictly before each forecast. On the seven known laws
(Fig.~\ref{fig:forecast}), unlabelled monitor features lowered the Brier score from
0.231 (time only) to 0.166, and released labels lowered it further to 0.149; the paired
gains were 0.0644 [0.0603, 0.0674] and 0.0169 [0.0160, 0.0179]. The released-label gain
shrinks with delay: Brier was 0.126 at delay 0, 0.155 at delay 5 and 0.166 at delay 20.
On a held-out conditional law with onsets balanced across four times, time alone still
gave the highest AUC (0.875). Monitor features improved the Brier score (0.238 to
0.188), but released labels gave no further gain ($-0.0067$ [$-0.0229$, 0.0086]), and
the current concern score ranked worse than chance (AUC 0.465). An earlier version of
this test with a shared onset gave a time-only AUC of 0.998; that design was
confounded and is reported as such.

\begin{figure}[t]
\centering
\resizebox{\columnwidth}{!}{% Data: data/forecast.dat <- results/phase3/future_harm_forecast_20261116_v1/evaluation_summary.json (overall)
% and results/phase3/future_harm_unseen_family_balanced_20261116_v1/metrics_overall_by_delay.csv (pooled rows).
\begin{tikzpicture}
\begin{groupplot}[
  group style={group size=2 by 1, horizontal sep=10mm},
  idm, width=0.58\columnwidth, height=4.3cm,
  ybar=0pt, /pgf/bar width=4.6pt, xmin=-0.6, xmax=4.6,
  xtick={0,...,4}, xticklabels={Prev.,Concern,Time,Monitor,Released},
  x tick label style={rotate=35, anchor=east, font=\scriptsize},
  xtick style={draw=none},
]
\nextgroupplot[title={Brier score (lower is better)}, ymin=0, ymax=0.7,
  legend style={legend columns=1, draw=cink!20}, legend to name=fcleg]
\addplot[fill=c1, draw=c1!80!black] table[x=idx, y=knownBrier] {data/forecast.dat};
\addlegendentry{known laws (C0--C6), 11,970 forecasts}
\addplot[fill=c5, draw=c5!80!black] table[x=idx, y=unseenBrier] {data/forecast.dat};
\addlegendentry{held-out law, balanced onset, 1,710 forecasts}
\nextgroupplot[title={ROC AUC (higher is better)}, ymin=0.4, ymax=0.95]
\addplot[fill=c1, draw=c1!80!black] table[x=idx, y=knownAUC] {data/forecast.dat};
\addplot[fill=c5, draw=c5!80!black] table[x=idx, y=unseenAUC] {data/forecast.dat};
\draw[cgray, dashed] (axis cs:-0.6,0.5) -- (axis cs:4.6,0.5);
\end{groupplot}
\node at ($(group c1r1.south)!0.5!(group c2r1.south)+(0,-16mm)$) {\pgfplotslegendfromname{fcleg}};
\end{tikzpicture}
}
\caption{Five-batch-ahead harm forecasting on held-out seeds: the seven known synthetic laws versus a held-out conditional law with balanced onset times. Prev.: development prevalence; Concern: current $\kappa_t$; Time: batch index only; Monitor: + unlabelled monitor features; Released: + released-label features. The dashed line marks chance AUC.}
\label{fig:forecast}
\end{figure}

On the natural streams, a five-cohort-ahead accuracy replay gave mixed results
(Table~\ref{tab:extforecast}). On ELEC2 the current monitor estimate (MAE 0.043) beat
the historical reference (0.094); on NOAA the historical reference (0.033) beat the
monitor (0.053); on UCI the last released accuracy (0.316 at delay 0) beat the frozen
monitor (0.585).

\subsection{Robustness}
\label{sec:res-robust}
\emph{Scale.} Repeating the core design with 1,500 seeds per task narrows the
intervals by an order of magnitude and leaves every conclusion unchanged
(Table~\ref{tab:largeseed}). Under C3 the calibrated MAE is 0.3028 [0.3024, 0.3032],
coverage is 0.310, and on average 16.08 harmful batches per stream go unflagged; C6
is matched at onset in 563 of 1,500 seeds (37.5\%). \emph{Learner.} The depth-6 tree
is less accurate than logistic regression on the linear task ($-0.0323$
[$-0.0343$, $-0.0301$] on C0) and much more accurate on stationary XOR ($+0.2168$
[0.1567, 0.2751]). Calibration lowers the tree's pooled MAE from 0.134 to 0.105, yet
the calibrated tree flags none of the 362 harmful batches under the XOR concept
complement or the 316 under increased label noise. \emph{Nuisance change.} A sinusoid
on an irrelevant feature raised the input-alarm rate by 0.744--0.768 across three
tasks, while every paired change in accuracy had an interval containing zero
(Table~\ref{tab:seasonal}). An alarm-driven policy would therefore spend most of its
attention on harmless batches.
